\section{Objective}\label{sec:intro}
Unsymmetric (Petrov-Galerkin) elements with parametric test and metric trial
functions reproduce every displacement field of total degree $\le p$ exactly on
arbitrarily distorted meshes, including meshes with $\det(\bm J)\le0$, where $\bm J$ is the
Jacobian matrix of the element map (\emph{distortion immunity})~\cite{Rajendran2003novel,Rajendran2010technique}.
Displacement fields of practical problems are, however, not polynomial.
\Citet{Rajendran2010technique} observed for the thick cylinder (Section~5.9 of that
paper, 8-node element) that the energy-norm convergence rate decreases from about
2 to 1.85, 1.42 and 1.21 as the distortion angle $\alpha$ (the angle by which the
internal mid-side nodes of every super-element are rotated, \cref{sec:mesh}) increases
from $0^\circ$ to $30^\circ$, $45^\circ$ and $60^\circ$, and concluded that the conditions C1--C3 of that paper, which are sufficient for
distortion immunity, do not guarantee optimal convergence for non-polynomial fields. This study revisits the
question for the unsymmetric spectral elements of the code base, $p=2,\dots,8$:
\begin{enumerate}
\item Is distortion resistance lost for non-polynomial fields?
\item If so, are optimal convergence rates nevertheless retained?
\item If both are lost, what is the practical benefit of unsymmetric elements?
\item Does a hybrid trial space, which uses metric (Cartesian) polynomials only for the
  complete space $P_p(\bx)$ and parametric functions for the remainder
  (\cref{sec:elements}), remove the deficiencies while keeping the polynomial exactness?
\end{enumerate}
The thick cylinder is complemented by manufactured solutions on the same geometry
and meshes: a polynomial field (degree 5), for which the element is immune for
$p\ge5$, and an analytic non-polynomial field. All statements below are based on the
computed data in \texttt{results/} (reproducible with the scripts listed in
\cref{sec:reproduce}); the conventional spectral element on the same meshes serves
as reference. Questions 1--3 concern the unsymmetric element \PM\ in its original form,
question 4 the hybrid element \PH.
